\documentclass[11pt,twocolumn]{article}
\usepackage[margin=2cm]{geometry}
\usepackage{amsmath,amssymb}
\usepackage{braket}
\usepackage{stix2}
\usepackage[hidelinks]{hyperref}

\begin{document}
\section*{Introduction}
Shadows, reflections and gravitational lensing in computer graphics, optics and astronomy are
commonly computed by ray tracing or path tracing, which assign each pixel the average of many
light paths and are the standard for physically based images. The mean over a pixel's rays is a
Monte Carlo integral. If the ray integrand can be evaluated coherently as an oracle, quantum amplitude
estimation (Brassard et al.\ 2002) estimates it to precision $\epsilon$ with $O(1/\epsilon)$ oracle
calls, against $O(1/\epsilon^2)$ independent samples (Montanaro 2015); applied to rendering, this is
quantum supersampling (Johnston 2016). Here the rays are precomputed, so the oracle encodes a finite
table of 1024 values per pixel, which a classical computer can also average exactly in 1024 reads.
The comparison with independent (IID) sampling therefore models a continuous integrand, and the
finite-table baseline is shown beside it. We examine its
query efficiency and its output on Kerr-Newman black holes with parameters motivated by the
supermassive black hole Sagittarius~A* (Sgr~A*).

\section*{Theory}

\subsection*{1. Spacetime and rays}
In Boyer-Lindquist coordinates and geometric units ($G = c = M = 1$), the Kerr-Newman metric
depends on the spin $a$ and charge $Q$ through
\begin{equation}
  \Sigma = r^2 + a^2\cos^2\theta, \qquad \Delta = r^2 - 2r + a^2 + Q^2,
\end{equation}
with outer horizon $r_+ = 1 + (1 - a^2 - Q^2)^{1/2}$. Photons are uncharged and therefore follow
null geodesics ($\mathrm{d}s^2 = 0$) of the metric. The Hamiltonian separates as
\begin{equation}
  2\Sigma H = \Delta\, p_r^2 + p_\theta^2 + \left(\frac{L}{\sin\theta} - a\sin\theta\right)^{2}
  - \frac{(r^2 + a^2 - aL)^2}{\Delta} = 0,
\end{equation}
with energy $E = 1$, axial angular momentum $L$ and the Carter constant conserved. Rays are
integrated backwards with fourth-order Runge-Kutta in Mino time $\tau$, defined by
$\mathrm{d}\lambda = \Sigma\,\mathrm{d}\tau$ with $\lambda$ the affine parameter, from a
zero-angular-momentum observer at $r = 40$ until they cross the horizon, escape to $r = 1000$ (star
field), or cross the disk plane. The traced shadow boundary agrees with the analytic photon-orbit
boundary to $7\times10^{-4}\,M$; the Carter constant drifts by less than $4\times10^{-7}$.

\subsection*{2. Disk emission}
We assume a geometrically thin, optically thin accretion disk between the innermost stable circular
orbit and $r = 10$. Gas follows prograde circular orbits with angular velocity
\begin{equation}
  \Omega = \frac{\sqrt{r - Q^2}}{r^2 + a\sqrt{r - Q^2}}.
\end{equation}
A photon of angular momentum $L$ emitted there is shifted in frequency by
$g = 1/\bigl(u^t(1 - \Omega L)\bigr)$. With $x = r_{\mathrm{ISCO}}/r$, the emitted flux follows the
thin-disk profile $F \propto x^3(1 - x^{1/2})$ and the temperature $T \propto F^{1/4}$; the observed
brightness is $g^4 F$ and the colour is the blackbody colour of $gT$. The temperature scale, lane
texture, filmic tone map and bloom filter are artistic and identical for every image shown.

\subsection*{3. Pixel integral and Monte Carlo}
A pixel's luminance is the mean of its $S = 1024$ ray values $I_j \in [0, 1]$,
$\mu = S^{-1}\sum_j I_j$. Averaging $N$ rays drawn at random gives an unbiased estimate $\hat\mu$
with root-mean-square error
\begin{equation}
  \epsilon_{\mathrm{MC}} = \sigma N^{-1/2},
\end{equation}
where $\sigma^2$ is the variance of the $I_j$. Smooth pixels have small $\sigma$ and converge
quickly; edges, the photon ring and stars have large $\sigma$.

\subsection*{4. Amplitude estimation}
With $n = \log_2 S$ index qubits and one ancilla, the circuit $A$ applies Hadamard gates to the
index register and a uniformly controlled rotation $R_y\bigl(2\arcsin I_j^{1/2}\bigr)$ to the
ancilla, so that
\begin{equation}
  A\ket{0} = \sqrt{1-\mu}\,\ket{\psi_0}\ket{0} + \sqrt{\mu}\,\ket{\psi_1}\ket{1},
  \qquad \mu = \sin^2\theta.
\end{equation}
The Grover operator $Q = A S_0 A^\dagger S_\chi$ acts as a rotation by $2\theta$ in the plane
spanned by $\ket{\psi_0}\ket{0}$ and $\ket{\psi_1}\ket{1}$, so after $k$ iterations the ancilla is
found in $\ket{1}$ with probability
\begin{equation}
  P_k(1) = \sin^2\!\bigl((2k+1)\theta\bigr).
\end{equation}
Maximum-likelihood amplitude estimation (Suzuki et al.\ 2020) measures $s = 32$ shots at each
$k \in \{0, 1, 2, 4, \dots, 128\}$, records the hits $h_k$, and maximizes
\begin{equation}
  \log\mathcal{L}(\theta) = \sum_k \bigl[h_k \log P_k + (s - h_k)\log(1 - P_k)\bigr].
\end{equation}
Counting each use of $A$ or $A^\dagger$ as one query, $N = s\sum_k (2k+1)$, and the error
approaches the Heisenberg scaling
\begin{equation}
  \epsilon_{\mathrm{AE}} \propto N^{-1},
\end{equation}
For this amplitude encoding, ideal shot probabilities depend on $\mu$, rather than separately
on the ray variance $\sigma^2$. The crossover against IID sampling depends on $\mu$, $\sigma$,
the schedule and estimator failures; $10^3$ is an empirical scale for selected pixels,
not a universal threshold. A finite table can instead be enumerated exactly.

\subsection*{5. Hybrid rule and error maps}
The hybrid spends the first 32 queries per pixel on Monte Carlo rays and computes
$\rho = \hat\sigma_{32}^2/\bigl(\hat\mu_{32}(1 - \hat\mu_{32})\bigr)$ from their sample mean
$\hat\mu_{32}$ and sample variance $\hat\sigma_{32}^2$. At budget $N$, a pixel with
$\rho > 300/N$ spends the remaining $N - 32$ queries on amplitude estimation: the schedule up to the
next $k$, with as many shots at that $k$ as fit (21 to 31). The pixel keeps its Monte Carlo estimate
if the result differs from $\hat\mu_{32}$ by more than four conservative pilot standard errors,
$\sqrt{\max(\hat\mu_{32}(1-\hat\mu_{32}), 1/32)/32}$. Other pixels continue Monte Carlo and reuse the 32 pilot rays. The advantage map marks the
first budget $N^*$ at which $\sqrt{\mu(1-\mu)/F(N^*)} < \sigma/\sqrt{N^*}$, with Fisher information
$F = s\sum_k (2k+1)^2$. The threshold 300 is inherited from the first version and was not
recalibrated on held-out data. The error display shows
$\log_{10}\lvert\hat\mu - \mu\rvert$ per pixel on a fixed scale from $10^{-5}$ to $10^{-1}$, so
passes and methods can be compared directly.

\subsection*{6. Limitations}
The target is the mean of quantized, tone-mapped grayscale ray values, not a calibrated radiance
prediction. Reference chroma is supplied classically to every displayed method. Rays and oracle
tables are precomputed; gate synthesis, data loading and ray tracing are excluded from the query
count. Full-image passes are computed from the exact noiseless output distribution with three random
seeds; Fig.~3 runs the circuits live. The finite table of stored rays is averaged exactly in 1024
classical reads, so these results are not evidence of a practical quantum speedup. The disk shading
is referred to an observer at infinity while the camera is at finite radius (factor 1.026 in $g$ for
the Sgr~A* camera at $L = 0$). The real charge of Sgr~A* is negligible (Zaja\v{c}ek et al.\ 2018),
so the charged preset is hypothetical.

\section*{References}
\begin{enumerate}\small
\item G. Brassard, P. H\o yer, M. Mosca and A. Tapp, Quantum amplitude amplification and estimation, \emph{Contemp.\ Math.} 305, 53 (2002).
\item Event Horizon Telescope Collaboration, First Sagittarius A* Event Horizon Telescope results. I., \emph{Astrophys.\ J.\ Lett.} 930, L12 (2022).
\item E. R. Johnston, Quantum supersampling, \emph{ACM SIGGRAPH 2016 Talks} (2016).
\item A. Montanaro, Quantum speedup of Monte Carlo methods, \emph{Proc.\ R. Soc.\ A} 471, 20150301 (2015).
\item Y. Suzuki et al., Amplitude estimation without phase estimation, \emph{Quantum Inf.\ Process.} 19, 75 (2020).
\item M. Zaja\v{c}ek, A. Tursunov, A. Eckart and S. Britzen, On the charge of the Galactic centre black hole, \emph{MNRAS} 480, 4408 (2018).
\end{enumerate}

\end{document}
